\section{System Architecture}
\label{sec:design_arch}

OpsMind AI is engineered following a decoupled, four-tier client-server architecture designed for modularity, low-latency execution, and operational reliability. The four distinct architectural layers comprise:

\begin{enumerate}
    \item \textbf{Presentation Layer (Frontend UI):} Built with React.js, Vite, and Tailwind CSS. It provides the user-facing Single Page Application (SPA), managing conversational interactions, file upload progress meters, document analytics dashboards, dark/light theme persistence, and interactive PDF source citation modals.
    \item \textbf{Application and Orchestration Layer (Backend Server):} Implemented using Node.js and Express.js. This layer handles HTTP routing, multi-part form data ingestion, PDF text parsing via \texttt{pdfjs-dist}, sliding-window text segmentation, in-process embedding coordination, and context prompt assembly.
    \item \textbf{Inference and Neural Embedding Layer:} Executes the local Transformer feature extraction pipeline (\texttt{all-MiniLM-L6-v2} via \texttt{@xenova/transformers}) to generate 384-dimensional dense vectors, and communicates asynchronously with Google's cloud-hosted Gemini 2.5 Flash API for context-grounded response generation.
    \item \textbf{Persistence and Vector Indexing Layer:} Backed by MongoDB Atlas. It stores chunk metadata, structural headers, character offsets, and high-dimensional vector embeddings, while executing native vector search aggregations via dedicated approximate nearest neighbor indices.
\end{enumerate}

Figure \ref{fig:arch_diagram} illustrates the comprehensive architectural topology and inter-component communication channels of OpsMind AI.

\begin{figure}[H]
\centering
\resizebox{0.95\textwidth}{!}{
\begin{tikzpicture}[
    comp/.style={rectangle, draw=blue!80!purple, fill=blue!12, text centered, rounded corners=4pt, minimum height=1.1cm, minimum width=4.0cm, text width=3.8cm, font=\small},
    layer_box/.style={rectangle, draw=blue!70, fill=blue!4, thick, rounded corners=6pt, inner sep=12pt},
    layer_title/.style={font=\bfseries\large, text=black, anchor=west},
    line/.style={draw=blue!80!black, ->, ultra thick, >=stealth}
]

    % Row 1: Presentation Layer Components
    \node [comp] (chat_ui) at (-4.5, 6.4) {Chat Assistant UI};
    \node [comp] (upload_ui) at (0, 6.4) {SOP Ingestion UI};
    \node [comp] (admin_ui) at (4.5, 6.4) {Admin Dashboard};
    
    % Row 1 Title & Background Box
    \node [layer_title] (pres_title) at (-6.6, 7.4) {Presentation Layer (React.js + Vite + Tailwind CSS)};
    \begin{scope}[on background layer]
        \node [layer_box, fit=(pres_title) (chat_ui) (admin_ui)] (pres_layer) {};
    \end{scope}

    % Row 2: Orchestration Layer Components
    \node [comp] (router) at (-4.5, 2.0) {API Routing \& Multer};
    \node [comp] (parser) at (0, 2.0) {PDF Text Extractor};
    \node [comp] (chunker) at (4.5, 2.0) {Sliding-Window Chunker};
    
    % Row 2 Title & Background Box
    \node [layer_title] (app_title) at (-6.6, 3.0) {Orchestration Layer (Node.js + Express.js)};
    \begin{scope}[on background layer]
        \node [layer_box, fit=(app_title) (router) (chunker)] (app_layer) {};
    \end{scope}

    % Row 3: Inference & Persistence Layer Components
    \node [comp] (embed) at (-4.5, -2.4) {all-MiniLM-L6-v2 (384-d)};
    \node [comp] (mongo) at (0, -2.4) {MongoDB Vector Search};
    \node [comp] (gemini) at (4.5, -2.4) {Google Gemini 2.5 Flash};
    
    % Row 3 Title & Background Box
    \node [layer_title] (infer_title) at (-6.6, -1.4) {AI Inference \& Persistence Layer};
    \begin{scope}[on background layer]
        \node [layer_box, fit=(infer_title) (embed) (gemini)] (infer_layer) {};
    \end{scope}

    % Connecting Arrows with Large Gap and Clear Labels
    \path [line] (pres_layer.south) -- node [right=10pt, font=\normalsize\bfseries, text=blue!80!black] {REST / JSON (Axios)} (app_layer.north);
    \path [line] (app_layer.south) -- node [right=10pt, font=\normalsize\bfseries, text=blue!80!black] {Embeddings, Aggregation \& Prompts} (infer_layer.north);

\end{tikzpicture}
}
\caption{Tiered System Architecture of OpsMind AI}
\label{fig:arch_diagram}
\end{figure}

\section{Data Flow and RAG Pipeline Architecture}
\label{sec:design_rag_pipeline}

The operational lifecycle of OpsMind AI is bifurcated into two coordinated execution pipelines: the \emph{Ingestion and Indexing Pipeline} and the \emph{Query and Retrieval Pipeline}.

\subsection{Document Ingestion and Indexing Pipeline}
The ingestion pipeline processes newly uploaded PDF manuals into searchable, high-dimensional vector representations through the following structured sequence of steps:

\begin{enumerate}
    \item \textbf{Step 1 (Upload Handling):} The client sends a multi-part form request containing the PDF binary. Multer intercepts the request, validates the MIME type, assigns a collision-resistant timestamped filename, and writes the file to the server storage directory.
    \item \textbf{Step 2 (Text Stream Extraction):} The PDF parsing engine (\texttt{pdfjs-dist}) reads the binary buffer, parses document cross-reference tables, iterates through each page sequentially, and extracts textual tokens into structured page arrays.
    \item \textbf{Step 3 (Sliding-Window Segmentation):} The extracted text for each page is divided into fixed-size chunks of 1,000 characters with a 100-character overlapping sliding window to ensure semantic continuity across split boundaries.
    \item \textbf{Step 4 (Section Header Detection):} Each chunk is evaluated against an operational domain dictionary to identify relevant operational topics and classify the segment under specific section headers.
    \item \textbf{Step 5 (Feature Extraction and Normalization):} The in-process ONNX neural pipeline transforms each text segment into a normalized 384-dimensional vector embedding.
    \item \textbf{Step 6 (Atomic Indexing):} Pre-existing chunks belonging to the same filename are purged to prevent duplicate records, and the newly generated chunk documents containing embeddings, metadata, and text are bulk-inserted into MongoDB Atlas.
\end{enumerate}

\subsection{Query and Retrieval Pipeline}
When an enterprise user submits an inquiry, the retrieval and answer generation pipeline executes as follows:

\begin{enumerate}
    \item \textbf{Step 1 (Query Vectorization):} The incoming natural language query string is passed to the embedding extractor, generating a normalized query vector $\mathbf{q} \in \mathbb{R}^{384}$.
    \item \textbf{Step 2 (Approximate Nearest Neighbor Search):} A MongoDB vector search aggregation pipeline evaluates the query vector $\mathbf{q}$ against the indexed corpus using cosine similarity, sampling 100 candidate vectors and returning the top 5 closest chunks.
    \item \textbf{Step 3 (Relevance Filtering):} The similarity score of the top-ranked chunk is checked against the threshold $\tau = 0.65$. If the score is below $\tau$, the pipeline terminates early and returns the standard refusal string.
    \item \textbf{Step 4 (Context Assembly):} If relevant chunks are found, the text contents of the top-$k$ chunks are concatenated into a formatted context block.
    \item \textbf{Step 5 (Generative Inference):} A strict anti-hallucination prompt containing the context and query is dispatched to the Gemini 2.5 Flash model.
    \item \textbf{Step 6 (Response and Citation Dispatch):} The synthesized answer, along with granular metadata for each supporting chunk (filename, page, section, chunk index, similarity score), is returned to the user interface.
\end{enumerate}

\section{Database Schema and Vector Index Design}
\label{sec:design_db_schema}

OpsMind AI utilizes a document-oriented database schema designed to seamlessly accommodate text chunks, structural metadata, and dense floating-point vector arrays.

\subsection{Chunk Collection Data Model}
Each record in the primary \texttt{chunks} collection adheres to the structural model outlined in Table \ref{tab:chunk_schema}.

\begin{table}[H]
\centering
\caption{Database Field Definitions for the Chunk Entity}
\label{tab:chunk_schema}
\resizebox{\textwidth}{!}{
\begin{tabular}{|l|l|p{8cm}|}
\hline
\textbf{Field Name} & \textbf{Data Type} & \textbf{Description} \\ \hline
\texttt{\_id} & \texttt{ObjectId} & Unique MongoDB document identifier \\ \hline
\texttt{fileName} & \texttt{String} & Original filename of the ingested SOP manual \\ \hline
\texttt{pdfPath} & \texttt{String} & Server-relative filesystem path to the uploaded PDF \\ \hline
\texttt{pageNumber} & \texttt{Number} & 1-indexed page number where the chunk originates \\ \hline
\texttt{section} & \texttt{String} & Classified operational section header or category \\ \hline
\texttt{chunkIndex} & \texttt{Number} & Monotonically increasing chunk sequence index \\ \hline
\texttt{text} & \texttt{String} & Raw extracted text content of the segment \\ \hline
\texttt{charCount} & \texttt{Number} & Total character count within the text segment \\ \hline
\texttt{embedding} & \texttt{Array<Number>} & 384-dimensional dense vector representation \\ \hline
\texttt{metadata} & \texttt{Object} & Boundary offsets (start and end character indices) \\ \hline
\end{tabular}
}
\end{table}

\subsection{Vector Search Index Specification}
To enable high-performance cosine similarity searches on the \texttt{embedding} field, an Atlas Vector Search index named \texttt{vector\_index} is configured with the parameters shown in Table \ref{tab:vector_index_spec}.

\begin{table}[H]
\centering
\caption{Vector Search Index Configuration}
\label{tab:vector_index_spec}
\resizebox{\textwidth}{!}{
\begin{tabular}{|l|l|p{7cm}|}
\hline
\textbf{Parameter} & \textbf{Setting} & \textbf{Technical Rationale} \\ \hline
\textbf{Index Name} & \texttt{vector\_index} & Identifier used in aggregation queries \\ \hline
\textbf{Field Path} & \texttt{embedding} & Target document field containing vector arrays \\ \hline
\textbf{Dimensions} & \texttt{384} & Native output dimension of \texttt{all-MiniLM-L6-v2} \\ \hline
\textbf{Similarity Metric} & \texttt{cosine} & Invariant to vector magnitude; evaluates angular orientation \\ \hline
\textbf{NumCandidates} & \texttt{100} & Number of nearest neighbor candidates evaluated during graph traversal \\ \hline
\textbf{Limit ($k$)} & \texttt{5} & Maximum top matching chunks returned for context compilation \\ \hline
\end{tabular}
}
\end{table}

\section{Text Segmentation and Domain Classification}
\label{sec:design_algorithms}

The text segmentation and domain classification procedure is used to divide the extracted SOP content into smaller, manageable sections and assign each section an appropriate operational domain. The procedure is designed to organize the document content while preserving the original meaning and semantics of the extracted text.

\subsection{Sliding-Window Segmentation}

The extracted text from each PDF page is processed using a sliding-window approach. A fixed window size of 1000 characters is used, with an overlap of 100 characters between consecutive chunks.

The segmentation process begins at the start of the extracted text. A portion of up to 1000 characters is selected and treated as an individual chunk. Each chunk is then associated with its corresponding page number and classified according to its operational section. After processing a chunk, the starting position is moved forward by 900 characters, resulting in a 100-character overlap with the previous chunk.

This overlapping strategy helps maintain contextual continuity between consecutive chunks, particularly when an important statement or procedure spans the boundary between two chunks. Each generated chunk contains its sequence number, text content, character count, page information, and identified operational section.

\subsection{Domain Section Classification}

After segmentation, each generated text chunk is analyzed to determine its relevant operational section. The classification procedure compares the content of the chunk with a predefined collection of domain-specific section patterns.

The pattern set includes sections such as:
\begin{itemize}
    \item Document Control
    \item Executive Summary
    \item Inbound Operations
    \item Inventory Audit Framework
    \item Outbound Operations
    \item Extreme Weather
    \item Other predefined operational sections contained within the SOP knowledge base.
\end{itemize}

The patterns are checked in a predefined order. When a matching section name is found within the chunk, that section is assigned as the classification label for the chunk. The classification process stops after the first matching pattern is identified.

If none of the predefined section patterns are detected in the chunk, the system assigns the default classification ``Unknown Section''.

This process allows the system to organize extracted SOP content according to operational domains, making the resulting chunks more meaningful for subsequent retrieval and question-answering operations.

\section{Context Grounding and Prompt Engineering}
\label{sec:design_prompts}

To guarantee factual adherence, prompt assembly enforces strict boundaries around the Large Language Model. The constructed prompt format isolates external knowledge:

\begin{equation}
\text{Prompt} = \mathcal{T}_{\text{system}} \oplus \mathcal{C}_{\text{retrieved}} \oplus \mathcal{Q}_{\text{user}}
\end{equation}

where $\mathcal{T}_{\text{system}}$ specifies the behavioral rules:
\begin{enumerate}
    \item The model must act exclusively as an authorized enterprise SOP assistant.
    \item The model is explicitly forbidden from utilizing external parametric knowledge or guessing.
    \item The answer must be derived solely from the provided context block $\mathcal{C}_{\text{retrieved}}$.
    \item If the factual answer is not explicitly stated within $\mathcal{C}_{\text{retrieved}}$, the model must reply with the exact string: \emph{``I don't know based on uploaded SOPs.''}
\end{enumerate}
